%%=============================================================================
%% CAPITULO 1 -- INTRODUCAO
%% Texto dummy do exercicio ia-6.1: dois paragrafos que marcam o lugar do
%% capitulo. Sera substituido pelo texto real da pesquisa.
%%=============================================================================

\section{Introdução}

O investidor pessoa física brasileiro que decide entre Tesouro Direto, CDB, LCI
e LCA no aplicativo do próprio banco raramente tem acesso a assessoria humana,
cujo custo não se paga em tíquetes pequenos. Recomendadores automatizados
prometem escala, e a IA generativa acrescenta a eles uma peça que a otimização
de carteira nunca teve: a justificativa da recomendação em linguagem natural,
com a qual o investidor pode dialogar.

Esta dissertação investiga se essa justificativa é, ao mesmo tempo, adequada ao
perfil declarado do investidor e fiel à ficha do produto recomendado. A
pergunta que conduz o trabalho é: em que medida recomendações de renda fixa
geradas por modelos de linguagem preservam a adequação ao perfil sem afirmar
nada que contrarie as características do produto? O fluxo analisado está
resumido na Figura~\ref{fig:fluxo}, no Capítulo~2.
